\section{EvoDiff: un modelo de difusión de secuencias de proteínas a escala evolutiva}

\subsection{Panorama del modelo y datos de entrenamiento}

EvoDiff es un \textbf{modelo generativo de proteínas basado en difusión discreta} desarrollado por el equipo de Amini en MSR, construido específicamente para el diseño de proteínas funcionales. El proyecto fue dirigido por Sarah Alamdari, con Amini y su colega Kevin Yang codirigiendo la orientación de la investigación.

\begin{importantbox}{Características centrales de EvoDiff}
\begin{itemize}
    \item \textbf{Tipo de modelo}: modelo generativo basado en difusión discreta (Discrete Diffusion)
    \item \textbf{Modalidad de datos}: secuencias de aminoácidos de proteínas (tokens discretos)
    \item \textbf{Datos de entrenamiento}: alrededor de 50 millones de secuencias de proteínas únicas, a escala evolutiva
    \item \textbf{Capacidad central}: generar nuevas secuencias de proteínas eliminando el ruido gradualmente a partir de un estado totalmente enmascarado
    \item \textbf{Capacidad extendida}: admite generación condicionada por función (motif scaffolding)
\end{itemize}
\end{importantbox}

«A escala evolutiva» significa que los datos de entrenamiento abarcan las secuencias de proteínas de distintos organismos a lo largo del árbol de la vida (tree of life), desde bacterias hasta humanos, e incluyen una representación de funciones biológicas extremadamente diversa. Esta amplia cobertura de datos permite que el modelo aprenda las regularidades generales de la distribución de las proteínas naturales.

\subsection{Dos modos de entrada de secuencia}

EvoDiff admite dos modos de entrada de secuencia, cada uno adecuado para un escenario de modelado distinto:

\textbf{Modo uno: secuencia única (Single Sequence)}

Se modela directamente una sola secuencia de proteína mediante difusión discreta. Es adecuado para el diseño de proteínas desde cero (de novo design) sin condicionamiento.

\textbf{Modo dos: Multiple Sequence Alignment (MSA)}

Aprovecha la información de contexto evolutivo. Para una secuencia objetivo dada, mediante una búsqueda de similitud de secuencias se encuentra un conjunto de \textbf{secuencias homólogas evolutivas relacionadas pero no idénticas}, estas secuencias se alinean y se disponen en una matriz (MSA) y, después, se ejecuta la difusión discreta sobre toda la matriz de MSA.

\begin{knowledgebox}{¿Por qué el MSA aporta información adicional?}
Multiple Sequence Alignment es un concepto clásico de la biología computacional. Las secuencias de proteínas relacionadas evolutivamente, aunque presentan mutaciones en ciertas posiciones, suelen mantenerse muy conservadas en las posiciones clave para la función. El patrón de conservación (conservation pattern) a nivel de columna en el MSA codifica de manera implícita la información crucial de \textbf{qué posiciones son fundamentales para la función}. El modo MSA de EvoDiff, al modelar de forma conjunta varias secuencias relacionadas, puede extraer un conocimiento previo más rico a partir de las señales evolutivas y así guiar el proceso de generación.
\end{knowledgebox}

\subsection{Visualización del proceso de generación}

El proceso de generación de EvoDiff puede entenderse de manera intuitiva así:

\begin{enumerate}
    \item \textbf{Inicialización}: una secuencia completamente en \texttt{[MASK]}, con la longitud de la proteína objetivo
    \item \textbf{Eliminación iterativa del ruido}: en cada paso, el modelo predice y rellena los aminoácidos de una parte de las posiciones enmascaradas
    \item \textbf{Formación de la secuencia}: a medida que avanzan los pasos, se rellenan cada vez más posiciones y la secuencia va tomando forma gradualmente
    \item \textbf{Finalización}: se rellenan todas las posiciones y se obtiene una secuencia de proteína completa
\end{enumerate}

Amini mostró en la clase una visualización dinámica: del lado izquierdo, el proceso de generación paso a paso del modelo de secuencias de EvoDiff; del lado derecho, sincronizada, la estructura tridimensional prevista que corresponde a la secuencia actual. Vale la pena enfatizar que \textbf{el aprendizaje y la generación de EvoDiff ocurren por completo en el espacio de secuencias}, sin utilizar ninguna información estructural; la estructura del lado derecho solo se usa para la visualización final.

\subsection{Resumen de la sección}

EvoDiff utiliza un marco de difusión discreta, entrenado sobre alrededor de 50 millones de secuencias de proteínas a escala evolutiva, y admite los dos modos de entrada de secuencia única y MSA. El modelo opera por completo en el espacio de secuencias, generando nuevas proteínas al eliminar el ruido gradualmente a partir de un estado totalmente enmascarado. El modo MSA, al incorporar información de contexto evolutivo, aporta un conocimiento previo más rico al proceso de generación.
